\section{Availability and Confluence under a Deployment}
\label{app:deployment}

Availability (Definition~\ref{def:partition-notions}) requires every participant
to be able to emit every output in its fiber.  A deployment restricts which rows
each participant may emit, assigning to each participant its own availability
relation where Theorem~\ref{thm:complete-cap} uses one for all.  This appendix restricts Availability accordingly and
characterizes the resulting compatibility condition.

\begin{definition}[Deployment, separable run family]
\label{def:deployment}
A \emph{deployment} over a finite participant set $P$ assigns to each $p\in P$
a set $A_p\subseteq S$, the rows $p$ may emit: $p$ performs an output event
$\Out{\gamma}{o}$ (Definition~\ref{def:system-output-event}) at current input
$i$ only if $(i,o)\in A_p$.  The \emph{emitted rows} of a run are the pairs
$(i,o)$ such that the run emits $o$ at current input $i$.  A \emph{partition} $\Pi$ of $P$ divides $P$ into
two or more nonempty components, as in a partitioned history
(Definition~\ref{def:partition-notions}).

A \emph{participant run family} assigns to each of finitely many participants
$p\in P$ an initial input and a canonical run whose emitted rows lie in $A_p$.
A participant whose run emits no output is \emph{silent}.  The family is
\emph{separable under $\Pi$} when either participants in at least two
components emit, or exactly one distinct row is emitted and some silent
participant lies in another component.
\end{definition}

\begin{definition}[Monotone across partitions]
\label{def:monotone-across-partitions}
Relative to a deployment, $\Spec$ is \emph{monotone across $\Pi$} when, for
every participant run family separable under $\Pi$ and every common \indesc $j$
of its participants' current inputs, some $u\in\fiber{j}$ \outrel every
emitted output.  It is \emph{monotone across partitions} when it is monotone
across every partition of $P$.
\end{definition}

This is the common-\indesc condition
(Definition~\ref{def:common-expansion-condition}) restricted to the families a
partition can separate, so Lifted Monotonicity implies monotonicity across
partitions.  Rows authorized only to $p$ can be emitted only by $p$'s run,
whose emitted outputs form a chain (each output event keeps only rows refining
it), so a family whose emitted rows are all authorized only to $p$ is
separable only when it emits a single row.  The converse can therefore fail.
For example, let $P=\{p,q\}$, let every fiber be $\{a,b\}$ ordered by equality,
and let $A_p=S$ and $A_q=\varnothing$, so that $p$ alone names a leader.  A run
of $p$ emits at most one distinct output and $q$ never emits, so every separable
family consists of one row $(i,a)$ or $(i,b)$ beside a silent $q$, and that
output lies in $\fiber{j}$ for every \indesc $j$ of $i$.  So $\Spec$ is
monotone across partitions, yet $(i,a)$ and $(i,b)$ have no common \outnoun and
Lifted Monotonicity fails.  The same happens whenever each authorized row has a
\outnoun at every \indesc of its input and every set of authorized rows
lacking a common \outnoun at some common \indesc is authorized only to a single
participant.  The product interface of Theorem~\ref{thm:product-frontier} and
the causal register of Appendix~\ref{app:registers} are further examples, with
no single authority, of monotonicity across partitions without Lifted
Monotonicity; they need not satisfy this condition, since a set of rows lacking
a common \outnoun may also contain rows of other participants.  If
$|P|\ge2$ and $A_p=S$ for every $p$, every pair of rows can be split across a
partition and every single row paired with a silent participant at the common
\indesc; since pairs of rows suffice for the common-\indesc condition
(Appendix~\ref{app:pairs-suffice}), monotonicity across partitions is then
Lifted Monotonicity.

\begin{definition}[Deployment Availability]
\label{def:deployment-availability}
Relative to a deployment, a system has \emph{Deployment Availability} when, at
every consistent cut of every partitioned history over $P$ it produces, each
participant $p$ in state $(i,R)$ can emit every $o\in R(i)$ with
$(i,o)\in A_p$ in some continuation, produced by the system and partitioned by
the same components, that adds only history events of $p$'s component and no
input event.
\end{definition}

As in Availability (Definition~\ref{def:partition-notions}), the output event
may apply commitments within the participant's component.  If $A_p=S$ for
every $p$, Deployment Availability is Availability for histories over $P$.

A deployment is \emph{single-authority} when $A_q=\varnothing$ for every
participant $q$ but one, as in the leader example above.  Every separable family
is then one row emitted by the single authority beside a silent participant, and
monotonicity across partitions says only that each such row has a \outnoun at
every \indesc of its input.  The formalism permits this case, as do
placement-sensitive models such as relational-transducer
CALM~\cite{ameloot2013relational,ameloot2015weaker}, and nothing below excludes
it.  It is degenerate for CAP, since a sole authority never needs to agree with
anyone, and we do not discuss it further.

\begin{proposition}[Partition compatibility under a deployment]
\label{thm:complete-cap-deployment}
Fix a deployment.  The following are equivalent.
\begin{enumerate}
  \item For every partition $\Pi$, every participant run family separable
        under $\Pi$ is confluent at every common \indesc $j$ of its
        participants' current inputs.
  \item $\Spec$ is monotone across partitions.
\end{enumerate}
\end{proposition}

\begin{proof}
Fix $\Pi$, a separable family, and $j$, and let $E$ be the set of emitted
outputs.  By Lemma~\ref{lem:reconciliation}, the family is confluent at $j$ if
and only if some $u\in\fiber{j}$ \outrel every output in $E$.  The two
conditions thus agree family by family.
\end{proof}

The proposition concerns canonical runs and does not use Deployment
Availability.  It applies to the runs of a system as well: by
Lemma~\ref{lem:canonical-equivalent}, each such run is observably equivalent to
a canonical run with the same emitted rows, which therefore lie in $A_p$.  So if
$\Spec$ is monotone across partitions, the outputs emitted by a separable
family of runs of a partitioned history, as Deployment Availability permits,
have a common \outnoun $u\in\fiber{j}$ at every common \indesc $j$.  It does not follow that the deployment authorizes any
participant to emit $u$; confluence through legal output events of the
deployment would additionally require that authority, a condition we do not
impose.  With $|P|\ge2$ and $A_p=S$ for every $p$, item~(2) is Lifted
Monotonicity, and the proposition is the step behind
Corollary~\ref{cor:complete-cap-S} stating that separable families are confluent
at every common \indesc exactly when Lifted Monotonicity holds.

\subsection{Registers}
\label{app:registers}

The Gilbert--Lynch witness, a replica cut off from a completed write that must
still answer a read, is recovered by Corollary~\ref{cor:gilbert-lynch}
(Appendix~\ref{app:cap-response}).  Here both writes are visible to both
replicas.

\paragraph{The causal register.}
This register records causal precedence rather than real-time precedence, and
its outputs record the write that supplies each read rather than the value
returned; it is otherwise the register of Appendix~\ref{app:cap-response}.  An
input is a finite operation history $h$: a set of reads and writes of one
register $x$ at replicas, with a strict partial order $\to$ of causal
precedence (program order at each replica, and from each write to the
operations a replica invokes after receiving it, closed transitively).  A set
of operations of $h'$ is \emph{causally closed} when it contains every
operation that causally precedes one of its members.  We set $h\InputOrd h'$
when the operations of $h$ form a causally closed subset of $h'$ and $\to$
restricted to them is the order of $h$.  Thus an \indesc never adds an
operation to the causal past of an existing one.

An output is a \emph{source map} $\rho$: a partial map assigning to some reads
of $h$ a write of $h$ or the initial value $v_0$, ordered by extension.  A
\emph{serialization} for replica $p$ is a total order of $p$'s operations and
all writes of $h$ that extends $\to$; it explains a read of $p$ with source
$w$ when $w$ is the last write before the read, or $v_0$ if there is none.
Under causal consistency~\cite{lamport1978time,ahamad1995causal},
$\rho\in\fiber{h}$ when each replica has a serialization explaining every read
of that replica in the domain of $\rho$.  Under linearizability, a single
total order of all operations of $h$ extending $\to$ must explain every read
in the domain~\cite{herlihy1990linearizability}; real-time edges beyond $\to$
play no role in the example below.

Two replicas start empty.  Participant $p$ completes $w_p(x,\dog)$ while $q$
concurrently completes $w_q(x,\catv)$; the replicas exchange both writes, a
partition separates them, and each serves a local read, $r_p$ at $p$ and $r_q$
at $q$.  Both writes causally precede both reads, and the writes are
concurrent.

Under linearizability, if $p$ reads \catv{} the order has
$w_p(x,\dog)$ before $w_q(x,\catv)$; if $q$ reads \dog{} it has the opposite.
No source map \outrel both, so the register is not monotone across the
partition $\{\{p\},\{q\}\}$ under any deployment authorizing these two reads.

Under causal consistency the two replicas may serialize the concurrent writes
differently, and one source map sends $r_p$ to $w_q(x,\catv)$ and $r_q$ to
$w_p(x,\dog)$.  Under the deployment in which each replica answers only its own
reads ($A_p$ holds the rows whose source map is defined only on reads of $p$),
the causal register is monotone across partitions, for every history.  Let
participant $p$'s run emit last the source map $\rho_p$ at input $i_p$, so that
$\rho_p$ \outrel every output $p$ emitted, and let $j$ be a common \indesc.
Take a serialization $L_p$ of $i_p$ explaining $\rho_p$, and extend it to $j$
by appending the new operations of $p$ and the new writes in any order
extending $\to$; since no new operation causally precedes an old one, the
result extends $\to$ and explains the same reads.  Using these serializations
for emitting replicas and arbitrary ones for the others, the union of the
$\rho_p$ lies in $\fiber{j}$ and \outrel every emitted output.  Separability is
not used.  Two answers to the same read remain incompatible, but they are
confined to one participant.  A commitment selecting a participant's response
is therefore local to its component.
